\documentclass[11pt]{article}
\usepackage[margin=1in]{geometry}
\usepackage[T1]{fontenc}
\usepackage[utf8]{inputenc}
\usepackage{lmodern}
\usepackage{microtype}
\usepackage{array}
\usepackage{booktabs}
\usepackage{enumitem}
\usepackage{hyperref}
\usepackage{xcolor}
\hypersetup{
  colorlinks=true,
  linkcolor=black,
  urlcolor=blue!50!black,
  pdftitle={Coding Scheme Dossier}
}
\setlist[itemize]{leftmargin=1.6em}
\setlist[description]{style=nextline,leftmargin=3.5cm,labelsep=0.6em}
\newcommand{\field}[1]{\nolinkurl{#1}}
\newcommand{\filepath}[1]{\nolinkurl{#1}}
\newcommand{\schemeheader}[3]{
  \begin{center}
    {\LARGE \textbf{#1}}\\[0.5em]
    {\large #2}\\[0.4em]
    {\normalsize \textit{#3}}
  \end{center}
  \vspace{0.8em}
}



\begin{document}
\schemeheader
  {Scheme 1}
  {Stage 1 Screening and Eligibility Decision Scheme}
  {Systematic Review on How the Biopsychosocial Model Frames Chronic Pain Research}

\section*{Purpose}
This scheme operationalizes title and abstract screening after search and deduplication. Its function is to decide whether a record should enter the review corpus for downstream coding, using a human-validatable rule set centered on biopsychosocial language, chronic pain relevance, review design, and population eligibility.

\section*{Operational Status}
\begin{itemize}
\item Workflow position: pre-extraction eligibility screen.
\item Operational mode: rule-based provisional machine assist with human validation.
\item Canonical implementation basis: executable screening rules plus protocol decision guidance.
\end{itemize}

\section*{Canonical Source Paths}
\begin{itemize}
\item \filepath{src/protocol/decision_rules/screening_rules.md}
\item \filepath{src/bps_review/screening/rules.py}
\item \filepath{src/review_stages/03_screening/README.md}
\item \filepath{src/protocol/osf/OSF_registration_HTBMFCPR.md}
\end{itemize}

\section*{Inputs Reviewed}
\begin{description}
\item[\field{title}] Record title.
\item[\field{abstract}] Record abstract text.
\item[\field{publication_types}] Publication type metadata used to distinguish reviews from protocols, commentaries, and primary studies.
\item[\field{language}] Language metadata used to enforce English-language eligibility.
\item[\field{publication_date} / \field{year}] Publication timing, checked against the operational search window.
\end{description}

\section*{Decision Categories}
\begin{description}
\item[\field{stage1_decision}] Implemented values are \texttt{include}, \texttt{exclude}, and \texttt{unclear}. The prose rules also allow a borderline \texttt{maybe} state; in the current executable implementation, ambiguous cases are operationalized as \texttt{unclear}.
\item[\field{stage1_reason}] Controlled reason attached to exclusions or unclear cases.
\item[\field{stage1_confidence}] Implemented values are \texttt{high}, \texttt{medium}, and \texttt{low}.
\item[\field{stage1_screened_by}] In current outputs this is recorded as \texttt{codex\_machine\_assist}.
\item[\field{stage1_screening_mode}] In current outputs this is recorded as \texttt{rule\_based\_provisional}.
\end{description}

\section*{Inclusion Logic}
\begin{itemize}
\item The record is a review article or review-like evidence synthesis.
\item The title or abstract explicitly contains a biopsychosocial term: \texttt{biopsychosocial}, \texttt{bio-psycho-social}, or \texttt{bio psycho social}.
\item The focus concerns chronic pain, persistent pain, or a chronic pain condition.
\item The population is adult or mixed-age rather than pediatric-only.
\item The record is within the operational search window and in English.
\end{itemize}

\section*{Exclusion Logic and Implemented Reason Labels}
\begin{description}
\item[\texttt{no biopsychosocial term in title/abstract}] No explicit biopsychosocial term is present in title or abstract.
\item[\texttt{outside operational search window}] Publication date falls before or after the operational search dates configured in \filepath{config/protocol.yaml}.
\item[\texttt{protocol}] Protocol papers are excluded.
\item[\texttt{commentary/editorial/letter}] Commentary-type publications are excluded.
\item[\texttt{animal/non-human focus}] Animal-only or non-human records are excluded unless a human focus is explicit.
\item[\texttt{pediatric-only focus}] Pediatric-only populations are excluded.
\item[\texttt{acute pain focus}] Acute pain records are excluded when chronicity is not also explicit.
\item[\texttt{chronic pain focus unclear}] Pain is present but chronic pain relevance is insufficiently clear.
\item[\texttt{non-English record}] Non-English records are excluded.
\item[\texttt{review status unclear}] Used when the record cannot be confidently identified as a review or evidence synthesis from title, abstract, or publication metadata.
\end{description}

\section*{Borderline Handling}
\begin{itemize}
\item Ambiguous review type, mixed acute/chronic populations, unclear age group, or unclear pain duration should not be over-excluded.
\item Borderline cases are logged with a rationale and lower confidence.
\item Musculoskeletal ambiguity is intentionally carried forward; Stage 3 full-text adjudication resolves that ambiguity rather than Stage 1 excluding too aggressively.
\end{itemize}

\section*{Primary Outputs Using This Scheme}
\begin{itemize}
\item \filepath{src/review_stages/03_screening/outputs/stage1_screening.csv}
\item \filepath{src/review_stages/03_screening/audit/stage1_screening_summary.csv}
\item \filepath{src/review_stages/03_screening/audit/reliability_report.csv}
\end{itemize}

\section*{Interpretive Note}
Although this scheme is not prompt-driven, it is part of the operational coding architecture of the review because it creates the standardized eligibility categories that determine which records can enter Stage 2 abstract coding and, later, Stage 3 full-text adjudication.

\end{document}

